% =============================================================================
% SECTION 3: EMERGENCE OF PHYSICAL QUANTITIES
% =============================================================================

\section{Emergence of Time, Mass, Charge, and Space}
\label{sec:emergence}

\subsection{The Time-Count Identity}
\label{subsec:time_count}

The partition count $M$ is defined as the cumulative number of completed
partition events in the history of a system. By Theorem~\ref{thm:floor}, $M$
is strictly monotone increasing. We now show that $M$ is proportional to
physical time.

\begin{theorem}[Time-Count Identity]
\label{thm:time_count}
For a system undergoing periodic partitioning with angular frequency $\omega$,
the rate of partition count accumulation is:
\begin{equation}
    \frac{dM}{dt} = \frac{\omega}{2\pi}
    \label{eq:time_count}
\end{equation}
Equivalently, the elapsed physical time between partition counts $M_1$ and
$M_2$ is:
\begin{equation}
    \Delta t = \frac{2\pi(M_2 - M_1)}{\omega} = \frac{\Delta M}{\dot{M}}
    \label{eq:time_from_count}
\end{equation}
\end{theorem}

\begin{proof}
Each partition event corresponds to one complete oscillatory cycle. One cycle
at frequency $\omega/(2\pi)$ Hz takes time $\tau_p = 2\pi/\omega$. Therefore
$\Delta M$ partition events take time $\Delta t = \Delta M \cdot \tau_p =
2\pi \Delta M / \omega$. Inverting gives $dM/dt = \omega/(2\pi)$.
\end{proof}

\begin{corollary}[Strict Monotonicity of Time]
\label{cor:time_mono}
Physical time, defined as $t = 2\pi M / \omega$, is strictly monotone
increasing. This follows directly from Theorem~\ref{thm:floor}: since $M$
cannot decrease, $t$ cannot decrease.
\end{corollary}

Corollary~\ref{cor:time_mono} is not a postulate about the arrow of time. It is
a theorem about the structure of partition events. The arrow of time does not
require an appeal to entropy, coarse-graining, or boundary conditions at the Big
Bang. It is a direct consequence of the irreducibility of the partition residue.

\subsection{Mass as Partition Inertia}
\label{subsec:mass}

\begin{definition}[Partition Inertia]
\label{def:inertia}
The \emph{partition inertia} of an item $\mathcal{I}$ is the quantity
$\mu(\mathcal{I})$ defined by:
\begin{equation}
    \mu = \alpha \cdot (m/z)
    \label{eq:inertia}
\end{equation}
where $m/z$ is the mass-to-charge ratio of $\mathcal{I}$, $\alpha$ is the
partition coupling constant, and the product $\mu$ determines how strongly the
item resists changes to its partition state.
\end{definition}

The physical mass $m$ of an item is its partition inertia in the limit of unit
charge ($z = 1$): $m = \mu/\alpha$. Mass is therefore not an intrinsic property
of matter but a measure of the item's resistance to partition state change ---
the cost, in partition count units, of transitioning from one partition state
$(n,\ell,m,s)$ to another.

\begin{theorem}[Mass-Frequency Relation]
\label{thm:mass_freq}
For an item with partition inertia $\mu$ undergoing oscillatory partition with
characteristic frequency $\omega_0$:
\begin{equation}
    m = \frac{\hbar \omega_0}{c^2}
    \label{eq:mass_freq}
\end{equation}
where $\hbar = E_{\mathrm{tick}}/(2\pi)$ is the reduced partition quantum and
$c = 2\pi$ rad/tick is the partition propagation rate.
\end{theorem}

\begin{proof}
The energy of one partition cycle is $E_{\mathrm{tick}} = \hbar \omega_0$.
By mass-energy equivalence (derived from the partition propagation rate $c$):
$E = mc^2$. Therefore $m = E/c^2 = \hbar\omega_0/c^2$.
\end{proof}

\subsection{Charge as Partition Asymmetry}
\label{subsec:charge}

The orientation index $m \in \{-\ell, \ldots, +\ell\}$ in the partition state
$(n,\ell,m,s)$ encodes the asymmetry of the boundary between $\mathcal{I}$ and
$\overline{\mathcal{I}}$. When $m = 0$, the boundary is symmetric: the item is
equally bounded on all sides. When $m \neq 0$, the boundary is asymmetric: the
item protrudes preferentially in one direction of the negation space.

\begin{definition}[Charge as Partition Orientation]
\label{def:charge}
The \emph{effective charge} of an item $\mathcal{I}$ in partition state
$(n,\ell,m,s)$ is:
\begin{equation}
    q = q_0 \cdot m
    \label{eq:charge_def}
\end{equation}
where $q_0$ is the elementary charge unit. Charge is the orientation of the
partition boundary --- its asymmetry relative to the global complement.
\end{definition}

\begin{theorem}[Charge Conservation]
\label{thm:charge_conserve}
For an isolated system of items $\{\mathcal{I}_k\}$ with total orientation
$\sum_k m_k = M_{\mathrm{orient}}$, the total charge is conserved:
\begin{equation}
    \sum_k q_k = q_0 M_{\mathrm{orient}} = \mathrm{const}
    \label{eq:charge_conserve}
\end{equation}
\end{theorem}

\begin{proof}
The orientation indices $\{m_k\}$ of the items in the system sum to
$M_{\mathrm{orient}}$, which is the net asymmetry of the system's boundary
with $\overline{\cup_k \mathcal{I}_k}$. Since the global complement is fixed
for an isolated system, and since the partition boundary between the system and
its complement has fixed total asymmetry, $M_{\mathrm{orient}}$ is constant.
\end{proof}

\begin{theorem}[Electromagnetic Interaction from Partition Orientation]
\label{thm:em_from_partition}
Two items $\mathcal{I}_1, \mathcal{I}_2$ with orientation indices $m_1, m_2$
interact through a potential:
\begin{equation}
    V(D) = \frac{q_0^2 m_1 m_2}{4\pi\epsilon_0} \cdot \frac{1}{r(D)}
    \label{eq:em_potential}
\end{equation}
where $r(D)$ is the spatial distance corresponding to partition depth $D$
(derived in Section~\ref{subsec:space}). Items with the same sign of $m$ have
parallel boundary orientations and repel (their local negation is less efficient
when they are close, since their boundaries point in the same direction and
compete). Items with opposite sign of $m$ have anti-parallel orientations and
attract (their local negation is more efficient when close, since each provides
what the other is negating).
\end{theorem}

\subsection{The S-Entropy Coordinates}
\label{subsec:sentropy}

For a partition state $(n,\ell,m,s)$, define the following normalised coordinates:

\begin{align}
    S_k &= \frac{n-1}{n_P - 1} \in [0,1]
    \label{eq:sk} \\
    S_t &= \frac{\ell}{n-1} \in [0,1] \quad (n > 1)
    \label{eq:st} \\
    S_e &= \frac{m + \ell}{2\ell} \in [0,1] \quad (\ell > 0)
    \label{eq:se}
\end{align}

where $n_P$ is the Planck depth (defined in Section~\ref{sec:composition}).
These three coordinates map every partition state to a point in $[0,1]^3$.

\begin{theorem}[S-Entropy as Partition Distance]
\label{thm:sentropy}
The S-entropy vector $\mathbf{S} = (S_k, S_t, S_e)$ measures the distance of
the partition state $(n,\ell,m,s)$ from the partition ground state $(1,0,0,s)$.
The Euclidean distance in $[0,1]^3$ is a well-defined metric on partition states.
\end{theorem}

\begin{proof}
The ground state $(1,0,0,s)$ has $S_k = 0$, $S_t = 0$, $S_e = 1/2$ (by
convention for the undefined case $\ell = 0$). Any other state has $S_k > 0$
and generally $S_t, S_e \neq 0, 1/2$. The Euclidean distance
$d = \sqrt{(S_k - 0)^2 + (S_t - 0)^2 + (S_e - 1/2)^2}$ is non-negative,
symmetric, and satisfies the triangle inequality, making it a metric.
\end{proof}

\subsection{Space as Emergent Partition Depth}
\label{subsec:space}

Space is not the container in which partition events occur. Space is the
coordinate representation of partition depth. We make this precise as follows.

\begin{definition}[Spatial Distance from Partition Depth]
\label{def:space}
The spatial distance $r(\mathcal{I}_1, \mathcal{I}_2)$ between two items is:
\begin{equation}
    r(\mathcal{I}_1, \mathcal{I}_2) = \lambda_P \cdot D(\mathcal{I}_1,
    \mathcal{I}_2)
    \label{eq:space}
\end{equation}
where $\lambda_P$ is the Planck length --- the spatial scale corresponding to
one unit of partition depth.
\end{definition}

This definition immediately gives:
\begin{enumerate}
    \item \textbf{Finite speed of signal propagation}: A signal propagating
          from $\mathcal{I}_1$ to $\mathcal{I}_2$ must traverse $D$ partition
          boundaries, each requiring time $\tau_p = 2\pi/\omega_0$ per boundary.
          Maximum speed: $c = \lambda_P/\tau_P = \lambda_P \omega_0/(2\pi)$.

    \item \textbf{Spatial dimensionality}: The three coordinates $(S_k, S_t,
          S_e)$ of S-entropy space define three independent directions of
          partition depth variation, giving rise to three spatial dimensions.

    \item \textbf{Curvature from partition density}: Regions with high
          partition count density (many items in close contact) have shorter
          effective $\lambda_P$ (more partition events per unit path length),
          manifesting as spacetime curvature in the coordinate representation.
\end{enumerate}

\subsection{Planck Depth and the Composition Limit}
\label{subsec:planck_depth}

\begin{definition}[Planck Depth]
\label{def:planck_depth}
For a system with $d$ distinguishable partition directions and oscillation
period $\tau_{\mathrm{osc}}$, the \emph{Planck depth} is:
\begin{equation}
    n_P = 1 + \left\lceil \log_{d+1} \left( \frac{\tau_{\mathrm{osc}}}{d \cdot
    \tau_P} \right) \right\rceil
    \label{eq:planck_depth}
\end{equation}
where $\tau_P = \sqrt{\hbar G / c^5}$ is the Planck time.
\end{definition}

For Cs-133 ($\tau_{\mathrm{osc}} = 1/9{,}192{,}631{,}770$ s, $d = 3$): $n_P = 56$.
This means the Cs-133 hyperfine transition operates at partition depth 56 --- it
requires 56 nested boundary layers to be well-defined as a distinct oscillatory
process. The partition count accumulated per second of Cs-133 oscillation is:
\begin{equation}
    \dot{M}_{\mathrm{Cs}} = 9{,}192{,}631{,}770 \;\mathrm{Hz}
    \label{eq:cs_mdot}
\end{equation}
which is by definition of the SI second.

\subsection{The Partition Lagrangian}
\label{subsec:lagrangian}

The dynamics of an item with partition inertia $\mu$ in a partition field are
governed by the Partition Lagrangian:

\begin{equation}
    \mathcal{L}_M = \frac{1}{2}\mu|\dot{x}|^2 + \mu\dot{x} \cdot A_M(x,t)
    - M(x,t)
    \label{eq:lagrangian}
\end{equation}

where $A_M(x,t)$ is the partition vector potential and $M(x,t)$ is the
partition count field. The Euler-Lagrange equations of
\eqref{eq:lagrangian} yield the equations of motion for the item in any
analyser geometry. We derive the four standard mass analyser equations as
consequences.

\begin{theorem}[Four Analyser Equations from Partition Lagrangian]
\label{thm:four_analysers}
The Euler-Lagrange equations of $\mathcal{L}_M$ in the appropriate boundary
conditions yield:
\begin{align}
    \text{TOF:} \quad T_{\mathrm{TOF}} &= L\sqrt{\frac{m}{2qV}}
    \label{eq:tof} \\
    \text{Orbitrap:} \quad \omega_z &= \sqrt{\frac{q\kappa}{m}}
    \label{eq:orbitrap} \\
    \text{FT-ICR:} \quad \omega_c &= \frac{qB}{m}
    \label{eq:fticr} \\
    \text{Quadrupole:} \quad a_z &= -2q_z \omega^2 / 4
    \label{eq:quad}
\end{align}
where $L$ is the TOF flight path length, $V$ is the accelerating voltage,
$\kappa$ is the Orbitrap geometry constant, $B$ is the FT-ICR magnetic field,
and $q_z$ is the Mathieu stability parameter.
\end{theorem}

\begin{proof}[Proof sketch]
Each analyser imposes a specific boundary condition on the partition field
$M(x,t)$. The TOF imposes a linear potential $M = qVx/L$; the Orbitrap imposes
a quadratic potential $M = \kappa x^2/2$; the FT-ICR imposes a vector potential
$A_M = B \times x/2$; the quadrupole imposes an oscillating potential
$M = (a - 2q\cos(2\Omega t))x^2/4$. Substituting each into the Euler-Lagrange
equations of \eqref{eq:lagrangian} and solving gives the respective frequency
or transit time equations above. Full derivations are given in
Appendix~\ref{app:analyser_derivations}.
\end{proof}

\subsection{The Partition Rates of the Four Analysers}
\label{subsec:partition_rates}

From Theorem~\ref{thm:time_count} and Theorem~\ref{thm:four_analysers}, the
partition rate $\dot{M} = \omega/(2\pi)$ for each analyser is:

\begin{align}
    \dot{M}_{\mathrm{Orb}}(m/z) &= \frac{1}{2\pi}\sqrt{\frac{q\kappa}{m}}
    \label{eq:mdot_orb} \\
    \dot{M}_{\mathrm{ICR}}(m/z) &= \frac{qB}{2\pi m}
    \label{eq:mdot_icr} \\
    \dot{M}_{\mathrm{TOF}}(m/z) &= \frac{1}{L}\sqrt{\frac{2qV}{m}} \cdot
    \frac{1}{2\pi}
    \label{eq:mdot_tof} \\
    \dot{M}_{\mathrm{quad}} &= \frac{q_u \Omega}{4\pi}
    \label{eq:mdot_quad}
\end{align}

Note that $\dot{M}_{\mathrm{quad}}$ is mass-independent at the stability apex
($q_u$ fixed), while all others scale as $(m/z)^{-1/2}$ (Orbitrap, TOF) or
$(m/z)^{-1}$ (FT-ICR). These scaling laws are direct predictions of the
Partition Lagrangian with zero free parameters.
